\documentclass[level=1]{kidscode}
\newcommand{\sticker}{\tikz[baseline=-0.35cm]\draw[kcgrey,dashed,rounded corners=3pt] (0,0) rectangle (1.5,-1.5);}
\newcommand{\lvboxes}{\small\tickbox\ Getting there\par\tickbox\ Got it\par\tickbox\ Expert}
%% \wkrow{week}{I can statements (\item ...)}{concept tags}
\newcommand{\ican}{\par\hangindent=1.2em\hangafter=1\noindent\faCheck\ }
\newcommand{\wkrow}[3]{\textbf{\large #1} &
  \vspace{-0.6\baselineskip}#2\par\smallskip
  {\footnotesize\color{kcgrey!70!black}#3} &
  \vspace{-0.6\baselineskip}\lvboxes & \sticker\\ \midrule}
%% breakable table (the dyslexia edition needs more than one page per half)
\makeatletter\ifkc@dyslexia\newcommand{\wkstretch}{1.25}\else\newcommand{\wkstretch}{1.9}\fi\makeatother
\newenvironment{weektable}{\par\noindent\renewcommand{\arraystretch}{\wkstretch}%
  \longtable{@{}>{\centering}p{1.15cm}>{\RaggedRight}p{\dimexpr\linewidth-5.95cm-6\tabcolsep}>{\RaggedRight}p{2.8cm}c@{}}
  \toprule
  \textbf{Week} & \textbf{I can\dots} & \textbf{How it went} & \textbf{Sticker}\\ \midrule\endhead}
  {\endlongtable\par}
%% badge: \badge{colour}{icon}{name}{how to earn it}{shape: circle|diamond|rectangle}
\newcommand{\badge}[5]{%
  \begin{tikzpicture}
    \node[#5,draw=#1!80!black,line width=1.6pt,fill=#1!12,minimum size=2.5cm,inner sep=1pt,
          align=center,font=\small\bfseries,text width=2.2cm] (b) {{\Large #2}\\[2pt] #3};
    \node[below=3pt of b,text width=3.8cm,align=center,font=\footnotesize] {#4\\[2pt] Date: \rule{1.8cm}{0.4pt}};
  \end{tikzpicture}}
\begin{document}
\kctitle{Progress Tracker}{One copy per student: what I can do, week by week}

\begin{tip}[How the tracker works]
At the end of each session (2 minutes, during the Wrap-up), read the \enquote{I can\dots} lines with the
student. The \textbf{student} ticks how it went, and you add a sticker or a stamp. Ticks are about
\emph{this week}: it is completely normal to tick \enquote{Getting there} for a new idea, because it comes
back later (look at the tags). Keep the trackers in a folder and bring them out at the showcase: students love seeing how far they have come.
\end{tip}

\begin{center}\small\renewcommand{\arraystretch}{1.3}
\begin{tabular}{@{}lll@{}}
\toprule
\textbf{Tag} & \textbf{Means} & \textbf{So\dots}\\
\midrule
\cvI & Introduced: brand new this week & \enquote{Getting there} is great.\\
\cvR & Revisited: we practise it again & Aim for \enquote{Got it}.\\
\cvE & Extended: we go further with it & Try for \enquote{Expert}.\\
\bottomrule
\end{tabular}
\end{center}

\begin{teachernote}[Levels]
\textbf{Getting there:} can do it with help or by copying an example. \textbf{Got it:} can do it alone.
\textbf{Expert:} can do it alone, explain it, and use it in a new way. The formal snapshots are the week-6
checkpoint and the week-12 exit assessment (see the Teacher Guide); this tracker is the everyday picture.
\end{teachernote}

\clearpage
\section{My progress: weeks 1--6 (Scratch)}
\noindent Nickname: \rule{6cm}{0.4pt}\par
\begin{weektable}
\wkrow{1}{\ican I can give exact step-by-step instructions in the right order.
  \ican I can make a sprite move when the green flag is clicked.}{sequence \cvI\ \ algorithm \cvI\ \ sprite, stage \cvI}
\wkrow{2}{\ican I can use a repeat loop to draw a shape with the pen.
  \ican I can work out the turn: $360 \div$ number of sides.}{loops \cvI\ \ costumes \cvI\ \ angles \cvI}
\wkrow{3}{\ican I can start code with events (flag, keys).
  \ican I can use if and if-else to make a sprite react.}{events \cvI\ \ if \cvI\ \ sequence \cvR\ \ loops \cvR}
\wkrow{4}{\ican I can use a variable to keep score.
  \ican I can use pick random for a surprise.}{variables \cvI\ \ random \cvI\ \ if \cvR\ \ events \cvR}
\wkrow{5}{\ican I can send and receive broadcast messages.
  \ican I can make a My Block with an input.}{broadcast \cvI\ \ My Blocks \cvI\ \ decomposition \cvI\ \ variables \cvR}
\wkrow{6}{\ican I can find and fix bugs in a Scratch project.
  \ican I can match Scratch blocks to Python lines.}{all Scratch ideas \cvR\ \ Scratch to Python \cvI}
\end{weektable}

\clearpage
\section{My progress: weeks 7--12 (Python)}
\noindent Nickname: \rule{6cm}{0.4pt}\par
\begin{weektable}
\wkrow{7}{\ican I can write and run a turtle program in Thonny.
  \ican I can find and fix typos and syntax errors.}{Python \cvI\ \ print, comments \cvI\ \ syntax errors \cvI\ \ sequence \cvR}
\wkrow{8}{\ican I can repeat with a \texttt{for} loop and \texttt{range}.
  \ican I can indent the lines that belong inside a loop.}{for/range \cvI\ \ indentation \cvI\ \ loops \cvE\ \ nested loop \cvI}
\wkrow{9}{\ican I can ask with \texttt{input} and change the answer to a number with \texttt{int}.
  \ican I can choose with \texttt{if}, \texttt{elif} and \texttt{else}.}{input \cvI\ \ comparisons \cvI\ \ if/elif/else \cvE\ \ variables \cvE}
\wkrow{10}{\ican I can make my own function with \texttt{def} and use it.
  \ican I can use \texttt{random} in Python and plan my project.}{functions \cvI\ \ My Blocks \cvR\ \ random \cvE\ \ decomposition \cvR}
\wkrow{11}{\ican I can build my project in small steps and test it.
  \ican I can give and use helpful feedback.}{small steps \cvI\ \ all ideas \cvR\ \ feedback \cvR}
\wkrow{12}{\ican I can show my project and explain my code.
  \ican I can say what I want to learn next.}{explaining code \cvE\ \ reflection \cvR}
\end{weektable}

\clearpage
\section{My final project milestones}
\noindent Nickname: \rule{6cm}{0.4pt}\qquad Project name: \rule{5cm}{0.4pt}\par\medskip
{\renewcommand{\arraystretch}{1.6}
\begin{tabularx}{\linewidth}{@{}>{\bfseries}p{3.4cm}X>{\centering}p{1.8cm}c@{}}
\toprule
Milestone & Done when\dots & Date & Sticker\\
\midrule
M1 Plan\newline{\normalfont\small end of week 10} & My plan sheet has a sketch, my questions and my function, and my teacher said OK. & & \sticker\\ \midrule
M2 Core drawing\newline{\normalfont\small week 11} & My function draws my main thing in at least two places. & & \sticker\\ \midrule
M3 Surprise and polish\newline{\normalfont\small week 11} & All six ingredients work and a partner playtested it. & & \sticker\\ \midrule
M4 Showcase\newline{\normalfont\small week 12} & I ran my project for visitors and explained one part. & & \sticker\\
\bottomrule
\end{tabularx}}

\section{My badges}
Your teacher gives a badge sticker (or a signature) when you earn it. Every badge can be earned by anyone, at any time.
\begin{center}
\badge{kcvermillion}{\faBug}{Bug Hunter}{Found and fixed 5 bugs, and said how.}{circle}\hfill
\badge{kcblue}{\faRedo}{Loop Legend}{Used a loop inside a loop.}{rectangle}\hfill
\badge{kcpurple}{\faUserFriends}{Kind Partner}{Great driver \emph{and} navigator; gave kind feedback.}{diamond}\hfill
\badge{kcgreen}{\faShield*}{Safety Star}{Chose a nickname and kept private things private.}{circle}
\par\medskip
\badge{kcorange}{\faTh}{Pixel Artist}{Made a pixel picture on grid paper (B1).}{rectangle}\hfill
\badge{kcsky}{\faMagic}{Function Wizard}{Made a function with two or more inputs.}{diamond}\hfill
\badge{kcyellow}{\faRecycle}{Remix Master}{Remixed a project and gave credit.}{circle}\hfill
\badge{kcblack}{\faStar}{Showcase Star}{Showed and explained my project.}{rectangle}
\end{center}

\clearpage
\section{Group overview (for the teacher)}
One line per student. Write the band: \textbf{J} Just starting, \textbf{G} Getting there, \textbf{O} Got it, \textbf{X} Expert.
Use the checkpoint and exit columns for the formal assessments; the week columns are your quick impression.
\begin{center}\small\renewcommand{\arraystretch}{2.3}
\begin{tabularx}{\linewidth}{|>{\RaggedRight}X|*{6}{c|}>{\columncolor{kcsky!15}}c|*{5}{c|}>{\columncolor{kcsky!15}}c|c|}
\hline
\textbf{Nickname} & 1 & 2 & 3 & 4 & 5 & 6 & \makecell{Check-\\point} & 7 & 8 & 9 & 10 & 11 & \makecell{Exit} & \makecell{Project\\/20}\\ \hline
 & & & & & & & & & & & & & & \\ \hline
 & & & & & & & & & & & & & & \\ \hline
 & & & & & & & & & & & & & & \\ \hline
 & & & & & & & & & & & & & & \\ \hline
\end{tabularx}
\end{center}
\textbf{Notes} (support or extension ideas, things to celebrate):
\answerlines{6}
\end{document}
